\section*{Заключение}
\addcontentsline{toc}{section}{Заключение}

В данной работе достигнута цель --- систематизированы физические механизмы полярных сияний, их типы и определяющие факторы.

Основные выводы:

\begin{enumerate}
    \item \textbf{Физика сияний} включает три ключевых этапа: пересоединение магнитных линий на границе магнитосферы, ускорение электронов (квазистатическое и альфвеновское), возбуждение атомов атмосферы с излучением фотонов.
    
    \item \textbf{Типы сияний} делятся на диффузные (равномерное свечение, альфвеновское ускорение) и дискретные (дуги, лучи, корона, квазистатическое ускорение). Цвет зависит от высоты и энергии частиц: красный ($\sim$200 км), зелёный (100--120 км), синий/фиолетовый (ниже 100 км).
    
    \item \textbf{Внешние факторы}: южная компонента IMF ($B_z<0$), скорость солнечного ветра, индекс Kp, фаза солнечного цикла (максимум), сезон (равноденствия).
    
    \item \textbf{Методы изучения}: наземные приборы (камеры, радары, магнитометры), космические миссии (THEMIS, Cluster, MMS, Swarm) и численные модели (MHD, кинетические, гибридные).
    
    \item \textbf{Перспективные направления}: гибридные модели, использование сияний для мониторинга атмосферы, машинное обучение. Неперспективны: поиск единого механизма ускорения, чистая статистика без физики, простое увеличение разрешения спутников.
\end{enumerate}

Полярные сияния перестали быть загадкой. Они стали важным индикатором космической погоды, что необходимо для безопасности спутников, линий электропередач и связи в высоких широтах~\cite{kivelson1995, pasachoff2013}.